% Round 6, study D: the orientation ball of rigid methane, alone. SO(3) drawn as the axis-angle ball (direction = axis,
% distance = angle, radius pi, the skin glued to its opposite point). The twelve rotations of T as red dots: the centre,
% two tetrahedra of third-turns (through the bonds at +2pi/3, through the faces at -2pi/3), three half-turns on the
% skin along the edge axes. One cell of the twelve (the Voronoi cell of the centre) as a wire octahedron; its
% vertices are the quarter-turns, the starred 4-cycle relabellings. Faces drawn flat (true faces bulge; edge
% midpoints sit at 70.5 degrees, vertices at 90). Coordinates from ball_check.py (scratch), camera az 18, el 22.
\documentclass[10pt,border=1mm]{standalone}
\usepackage[T1]{fontenc}
\usepackage{figurestyle}
\input{primitives.tex}
% emitted by ball_check.py (scratch); az 18, el 22, 0.8 cm per radian
\def\ballR{2.5133}\def\ballE{0.3746}
\def\Hback{c/3,d/4}\def\Hfront{a/1,b/2}
\def\Hxa{0.2446}\def\Hya{0.1734}\def\Hda{0.5879}
\def\Cxa{0.6211}\def\Cya{0.4403}\def\Cda{1.4926}
\def\Dxa{-0.6211}\def\Dya{-0.4403}\def\Dda{-1.4926}
\def\Hxb{-0.4802}\def\Hyb{-0.4450}\def\Hdb{0.0841}
\def\Cxb{-1.2189}\def\Cyb{-1.1296}\def\Cdb{0.2135}
\def\Dxb{1.2189}\def\Dyb{1.1296}\def\Ddb{-0.2135}
\def\Hxc{0.4802}\def\Hyc{-0.2617}\def\Hdc{-0.3696}
\def\Cxc{1.2189}\def\Cyc{-0.6643}\def\Cdc{-0.9382}
\def\Dxc{-1.2189}\def\Dyc{0.6643}\def\Ddc{0.9382}
\def\Hxd{-0.2446}\def\Hyd{0.5332}\def\Hdd{-0.3024}
\def\Cxd{-0.6211}\def\Cyd{1.3535}\def\Cdd{-0.7678}
\def\Dxd{0.6211}\def\Dyd{-1.3535}\def\Ddd{0.7678}
\def\Sxpx{-0.7766}\def\Sxpy{-0.8954}\def\Sxpd{2.2162}
\def\Vxpx{-0.3883}\def\Vxpy{-0.4477}\def\Vxpd{1.1081}
\def\Sxmx{0.7766}\def\Sxmy{0.8954}\def\Sxmd{-2.2162}
\def\Vxmx{0.3883}\def\Vxmy{0.4477}\def\Vxmd{-1.1081}
\def\Sypx{2.3903}\def\Sypy{-0.2909}\def\Sypd{0.7201}
\def\Vypx{1.1951}\def\Vypy{-0.1455}\def\Vypd{0.3600}
\def\Symx{-2.3903}\def\Symy{0.2909}\def\Symd{-0.7201}
\def\Vymx{-1.1951}\def\Vymy{0.1455}\def\Vymd{-0.3600}
\def\Szpx{0.0000}\def\Szpy{2.3303}\def\Szpd{0.9415}
\def\Vzpx{0.0000}\def\Vzpy{1.1651}\def\Vzpd{0.4707}
\def\Szmx{0.0000}\def\Szmy{-2.3303}\def\Szmd{-0.9415}
\def\Vzmx{0.0000}\def\Vzmy{-1.1651}\def\Vzmd{-0.4707}
\def\cellVisible{Vxp/Vym,Vxp/Vyp,Vxp/Vzm,Vxp/Vzp,Vym/Vzm,Vym/Vzp,Vyp/Vzm,Vyp/Vzp}
\def\cellHidden{Vxm/Vym,Vxm/Vyp,Vxm/Vzm,Vxm/Vzp}
\def\cellFront{Vxp/Vyp/Vzp,Vxp/Vyp/Vzm,Vxp/Vym/Vzp,Vxp/Vym/Vzm}

\begin{document}
\begin{tikzpicture}[plate]
\path[use as bounding box] (0,0) rectangle (15.2,9.3);
\begin{scope}[shift={(7.4,4.75)}]
  % the ball: a glass sphere, teal because it takes over the core sheet's role
  \fill[ColBand] (0,0) circle[radius=\ballR];
  \draw[fine] (0,0) circle[radius=\ballR];
  \draw[hidden] ({-\ballR},0) arc[start angle=180,end angle=0,x radius=\ballR,y radius={\ballE*\ballR}];
  \draw[fine] ({-\ballR},0) arc[start angle=180,end angle=360,x radius=\ballR,y radius={\ballE*\ballR}];
  % two axes through the centre: the C2 axis z (skin to skin, through two cell vertices) and the C3 axis of bond 1
  \draw[hidden] (\Szmx,\Szmy)--(\Szpx,\Szpy);
  \draw[hidden] (\Dxa,\Dya)--(\Cxa,\Cya);
  % the cell: hidden edges, then visible edges
  \foreach \p/\q in \cellHidden {\draw[hidden] (\csname \p x\endcsname,\csname \p y\endcsname)--(\csname \q x\endcsname,\csname \q y\endcsname);}
  \foreach \p/\q in \cellVisible {\draw[fine] (\csname \p x\endcsname,\csname \p y\endcsname)--(\csname \q x\endcsname,\csname \q y\endcsname);}
  % the reference at the centre: back hydrogens, carbon, front hydrogens
  \foreach \k/\n in \Hback {\draw[molbond,line width=.5pt] (0,0)--(\csname Hx\k\endcsname,\csname Hy\k\endcsname); \node[atomH,minimum size=2.0mm,line width=.4pt,font=\fontsize{5}{5}\selectfont] at (\csname Hx\k\endcsname,\csname Hy\k\endcsname) {\n};}
  \node[atomC,minimum size=2.3mm] at (0,0) {};
  \foreach \k/\n in \Hfront {\draw[molbond,line width=.5pt] (0,0)--(\csname Hx\k\endcsname,\csname Hy\k\endcsname); \node[atomH,minimum size=2.0mm,line width=.4pt,font=\fontsize{5}{5}\selectfont] at (\csname Hx\k\endcsname,\csname Hy\k\endcsname) {\n};}
  % the twelve rotations of T: the centre (under the carbon), eight third-turns, three half-turns on the skin (each twice)
  \foreach \k in {a,b,c,d} {\fill[ColDot] (\csname Cx\k\endcsname,\csname Cy\k\endcsname) circle[radius=.075]; \fill[ColDot] (\csname Dx\k\endcsname,\csname Dy\k\endcsname) circle[radius=.075];}
  \foreach \a in {x,y,z} {\fill[ColDot] (\csname S\a px\endcsname,\csname S\a py\endcsname) circle[radius=.075]; \fill[ColDot!45] (\csname S\a mx\endcsname,\csname S\a my\endcsname) circle[radius=.075];}
  % the cell vertices: the quarter-turns, not in T (ink rings)
  \foreach \a in {x,y,z} {\foreach \t in {p,m} {\draw[fine,fill=white] (\csname V\a\t x\endcsname,\csname V\a\t y\endcsname) circle[radius=.05];}}
  % labels, with leaders that start on the object and stop short of the label
  \node[sub,anchor=south west,align=left] at ({-\ballR-.6},{\ballR+.3}) {the orientation ball $SO(3)$: direction the axis, distance the angle;\\the skin is the half-turns, each point glued to its opposite};
  \draw[leader] (-.14,.06)--({-\ballR-.35},.75); \node[sub,anchor=east,align=right] at ({-\ballR-.4},.75) {the reference,\\at the centre};
  \draw[leader] ($(\Cxa,\Cya)+(.08,.03)$)--({\ballR+.35},1.55); \node[sub,anchor=west,align=left] at ({\ballR+.4},1.55) {a third-turn about the bond to H$_1$:\\back with $2,3,4$ cycled; one per bond,\\a tetrahedron of four};
  \draw[leader] ($(\Dxa,\Dya)+(-.06,-.06)$)--({-\ballR-.35},-1.45); \node[sub,anchor=east,align=right] at ({-\ballR-.4},-1.45) {the opposite turn about\\the same bond; the four\\make the second tetrahedron};
  \draw[leader] ($(\Szpx,\Szpy)+(.08,.0)$)--({\ballR+.35},.35); \node[sub,anchor=west,align=left] at ({\ballR+.4},.35) {a half-turn about an edge axis,\\on the skin; three of them,\\each seen twice};
  \draw[leader] ($(\Vxpx,\Vxpy)+(.0,-.07)$)--(.7,{-\ballR-.5}); \node[sub,anchor=north,align=center] at (0,{-\ballR-.6}) {one cell of the twelve, Albert et al.'s octahedral space:\\the orientations nearer the centre than any other red dot; its vertices are the quarter-turns, the starred relabellings;\\each face is glued to the opposite face by a third-turn (faces drawn flat here; they bulge)};
\end{scope}
\end{tikzpicture}
\end{document}
